%% —-----------------------------------------------------------
%% sbcreviews-pt.tex
%% Template para SBC Reviews 
%% (https://journals-sol.sbc.org.br/index.php/reviews/) 
%%
%% Licença CC BY 4.0 (https://creativecommons.org/licenses/by/4.0/)
%% Sociedade Brasileira de Computação.

\documentclass[portuguese]{sbcreviews-2025}%
\usepackage[misc,geometry]{ifsym} 

\usepackage{aas_macros}
\usepackage[bottom]{footmisc}
\usepackage{tabularray}
\usepackage{afterpage}
\usepackage{url}
\usepackage{pifont}

\usepackage{tikz}
\usetikzlibrary{positioning,fit,backgrounds}
\tikzset{
    ator/.style={
        draw,
        rectangle,
        rounded corners,
        align=center,
        minimum width=2.8cm,
        minimum height=0.9cm,
        font=\scriptsize\bfseries,
        fill=gray!10
    },
    casodeuso/.style={
        draw,
        ellipse,
        align=center,
        minimum width=3.7cm,
        minimum height=0.9cm,
        font=\scriptsize,
        fill=white
    },
    associacao/.style={
        draw,
        line width=0.4pt
    }
}

\setcitestyle{square}

\definecolor{engtitle}{rgb}{0.5,0.5,0.5}
\definecolor{orcidlogo}{rgb}{0.37,0.48,0.13}
\definecolor{unilogo}{rgb}{0.16, 0.26, 0.58}
\definecolor{maillogo}{rgb}{0.58, 0.16, 0.26}
\definecolor{darkblue}{rgb}{0.0,0.0,0.0}
\hypersetup{colorlinks,breaklinks,
            linkcolor=darkblue,urlcolor=darkblue,
            anchorcolor=darkblue,citecolor=darkblue}

\jid{ROCS}
\jtitle{SBC Computing Reviews, 202X, XX:1}
\issn{2966-3938}
\doi{10.5753/rocs.202X.XXXXXX}
\copyrightstatement{This work is licensed under a Creative Commons Attribution 4.0 International License}
\jyear{2026}

\category{Trabalho Prático}
\title[DevMatch]{DevMatch: Especificação e Modelagem de uma Plataforma de Marketplace para Contratação de Profissionais de TI Freelancers}
\engtitle{\textcolor{engtitle}{DevMatch: Specification and Modeling of a Marketplace Platform for Hiring Freelance IT Professionals}}

%THE ORCID IS MANDATORY FOR EACH AUTHOR
\author[Andrade et al. 2026]{
\affil{\textbf{Antônio Henrique Rabelo Andrade}~\orcidlink{0000-0000-0000-0000}~[{PUC Minas}~|\href{mailto:...}{~{\textit{ahrandrade@sga.pucminas.br}}}~]}

\affil{\textbf{João Pedro Bomfim Victor}~\orcidlink{0000-0000-0000-0000}~[{PUC Minas}~|\href{mailto:...}{~{\textit{joao.bomfim@sga.pucminas.br}}}~]}

\affil{\textbf{Júlia de Souza Ventura}~\orcidlink{0000-0000-0000-0000}~[{PUC Minas}~|\href{mailto:...}{~{\textit{jsventura@sga.pucminas.br}}}~]}

\affil{\textbf{Pedro Henrique Germano de Melo Santos}~\orcidlink{0000-0000-0000-0000}~[{PUC Minas}~|\href{mailto:...}{~{\textit{pedro.germano@sga.pucminas.br}}}~]}

\affil{\textbf{Yuri Cardoso Viana}~\orcidlink{0000-0000-0000-0000}~[{PUC Minas}~|\href{mailto:...}{~{\textit{ycviana@sga.pucminas.br}}}~]}
}

\begin{document}

\begin{frontmatter}

\maketitle

\begin{mail}
Curso de Engenharia de Software, Pontifícia Universidade Católica de Minas Gerais (PUC Minas), Belo Horizonte, MG, Brasil.
\end{mail}

\begin{abstract-pt}
% Resumo em português (500–750 palavras) — Julia, a ser escrito por último
\end{abstract-pt}

\begin{abstract-en}
% Abstract em inglês — Julia, a ser escrito por último
\end{abstract-en}

\begin{pchaves}
Marketplace de Freelancers, Engenharia de Software, Modelagem de Sistemas, UML
\end{pchaves}

\begin{keywords}
Freelancer Marketplace, Software Engineering, System Modeling, UML
\end{keywords}

\end{frontmatter}

\section{Introdução}
\label{sec:intro}
% Antônio — contextualização do problema, apresentação do DevMatch, objetivos e organização do artigo

\section{Trabalhos Relacionados}
\label{sec:relacionados}
% Julia — análise de sistemas similares (Workana, Upwork, 99Freelas), lacunas e diferenciais do DevMatch

\section{Modelagem de Requisitos}
\label{sec:requisitos}

\subsection{Descrição de Atores}
\label{sec:atores}
% Antônio

\subsection{Modelo de Casos de Uso e Histórias de Usuário}
\label{sec:casosdeuso}
% Antônio

\subsection{Diagrama de Sequência do Sistema e Contratos de Operação}
\label{sec:seqsistema}
% Pedro Henrique — entrega 13/11

\section{Modelagem de Projeto}
\label{sec:projeto}

\subsection{Arquitetura do Sistema}
\label{sec:arquitetura}
% João Pedro

\subsection{Diagramas de Componentes e de Implantação}
\label{sec:componentes}
% João Pedro

\subsection{Diagrama de Classes}
\label{sec:classes}
% Yuri

\subsection{Diagramas de Sequência}
\label{sec:sequencia}
% Pedro Henrique — entrega 13/11

\subsection{Diagramas de Comunicação}
\label{sec:comunicacao}
% Pedro Henrique — entrega 11/12

\subsection{Diagramas de Estados}
\label{sec:estados}
% Pedro Henrique — entrega 11/12

\section{Modelagem de Dados}
\label{sec:dados}
% Julia — modelo conceitual, lógico e físico do banco de dados

\section{Conclusão}
\label{sec:conclusao}
% A ser escrito por último, com contribuição de todos

\begin{declarations}

\begin{acknowledgements}
% Opcional
\end{acknowledgements}

\begin{funding}
% Opcional
\end{funding}

\begin{contributions}
% Obrigatória — descrever a contribuição de cada integrante (Taxonomia CRediT)
\end{contributions}

\begin{interests}
Os autores declaram que não têm nenhum conflito de interesses.
\end{interests}

\begin{materials}
% Obrigatória — disponibilidade dos materiais (diagramas, modelos) usados no trabalho
\end{materials}

\begin{aitools}
% Obrigatória — declarar ferramentas de IA generativa utilizadas (ex.: Claude), onde foram empregadas, e que os autores mantêm total responsabilidade pelo conteúdo
\end{aitools}

\begin{furtherinformation}
Não se aplica aprovação por comitê de ética a este trabalho.
\end{furtherinformation}

\end{declarations}

\bibliographystyle{apalike-sol}
\bibliography{refs}

\end{document}